\chapter{Experimentos e Resultados}\label{chap:experiments}
\section{Experimentos}\label{sec:experiments}

Uma vez integrado ao Ket, o novo simulador pode ser utilizado naturalmente em conjunto à infraestrutura já existe da biblioteca. Ou seja, algoritmos quânticos já implementados podem suavemente ser executados tanto pelos simuladores já existentes quanto pelo novo simulador. Isso permite que o desempenho do simulador seja medido e comparado com os simuladores já existentes dentro do ambiente no qual serão utilizados.

Assim, o desempenho do novo simulador foi comparado com dois simuladores já existentes na plataforma Ket: o simulador denso e o simulador esparso. Essas comparações foram feitas utilizando diferentes algoritmos quânticos conhecidos na literatura: algoritmo de Grover, FALQON e algoritmo de estimação de fase quântica. Cada algoritmo foi simulado utilizando diferentes quantidades de qubits para demonstrar como a complexidade de cada algoritmo afeta o desempenho dos simuladores.

Todos os experimentos foram realizados em um mesmo computador. A seguir estão as especificações de hardware e software do computador utilizado para a execução dos experimentos:

\begin{itemize}
  \item Processador: Intel 7 240H (16) @ 5.0GHz
  \item Memória RAM: 32GB 1300MHz
  \item Sistema Operacional: openSUSE Leap 16.0 x86\_64
  \item Kernel: Linux 6.12.0
  \item Compilador: rustc 1.98.0
  \item Python: 3.13.5
  \item Ket: 0.10.4.1
\end{itemize}

\section{Simuladores}\label{sec:simulators}

\subsection{Simulação Densa}\label{sub:dense}

Na simulação densa, um circuito de $n$ qubits é simulado utilizando um array de números complexos com $2^n$ posições. Nesse array, cada posição representa a amplitude de um dos $2^n$ possíveis estados do circuito. Assim, cada porta é aplicada individualmente sobre as posições do array que estejam atreladas ao qubit afetado pela porta.

Por exemplo, um circuito de 2 qubits no estado $\frac{1}{\sqrt{2}}(\ket{00} + \ket{11})$ é representado pelo array a abaixo:

\[
[\frac{1}{\sqrt{2}}, 0, 0, \frac{1}{\sqrt{2}}]
\]

\subsection{Simulação Esparsa}\label{sub:sparse}

Já a simulação esparsa se baseia em tabelas que armazenam apenas as amplitudes não nulas do circuito. No Ket, existem duas implementações desse método: uma utilizando tabelas hash e outra utilizando vetores. Ao utilizar tabelas hash, cria-se pares de chave e valor, sendo a chave a string de bits que representa um determinado estado e o valor a sua amplitude. 

Por exemplo, o estado $\frac{1}{\sqrt{2}}(\ket{00} + \ket{11})$ seria representado pela tabela.

No Ket, esse simulador é implementado de forma a permitir a paralelização do cálculo de aplicação das matrizes sobre o vetor de estado. Isto é, o vetor de estados pode ser divido em vários trechos de forma que diferentes trechos possam ser processados simultaneamente por diferentes núcleos de um mesmo computador.

\begin{table}[ht]
  \centering
  \begin{tabular}{|c|c|}
    \hline
    Estado & Amplitude \\
    \hline
    00 & $\frac{1}{\sqrt{2}}$ \\
    11 & $\frac{1}{\sqrt{2}}$ \\
    \hline
  \end{tabular}
  \caption{Exemplo de representação esparsa de um estado quântico.}
  \label{tab:sparse-state-example}
\end{table}

\section{Algoritmos Quânticos}\label{sec:quantum_algorithms}

Dada a natureza variável do desempenho do QMDD, foram utilizados diferentes algoritmos quânticos para identificar os cenários nos quais o novo simulador apresenta vantagens ou desvantagens em relação aos demais. As subseções seguintes descrevem esses algoritmos e destacam suas características pertinentes ao desempenho da simulação.

\subsection{Algoritmo de Grover}\label{sub:grover}

Em 1996, Grover propôs um algoritmo quântico para busca em conjuntos desordenados. Esse algoritmo é composto principalmente de portas X, H e CZ (portas Z controladas). A complexidade desse algoritmo surge, inicialmente, da necessidade de repetir a mesma operação múltiplas vezes. \cite{grover}

Porém, para a simulação por QMDD, o fato de todos os qubits do circuito agirem como qubits de controle para as portas Z faz com que o tempo de simulação cresça exponencialmente em relação ao número de qubits do circuito simulado.

\subsection{Algoritmo FALQON}\label{sub:falqon}

{\itshape Feedback-based ALgorithm for Quantum OptimizatioN} (FALQON) é um algoritmo de otimização generalizado cuja proposta é eliminar a parte clássica do processo iterativo dos algoritmos de otimização quânticos até então propostos. \cite{magann2022falqon}

Por ser um algoritmo de otimização, o FALQON requer um problema de otimização a ser resolvido em cada experimento. Para isso, foi utilizado o exemplo disponível na página de tutoriais do Ket: o problema max-cut \cite{ket2021falqon}. Esse problema envolve grafos e sua complexidade é exponencial em relação ao número de arestas presentes no grafo, tornando-o adequado à avaliação do algoritmo.

O FALQON também utiliza portas controladas, porém em quantidade menor que o algoritmo de Grover. Apesar de manter complexidade exponencial, esse algoritmo apresenta custo reduzido ao ser simulado por QMDDs em comparação ao algoritmo de Grover.

\subsection{Algoritmo de Estimação de Fase Quântica}\label{sub:qpe}

Estimação de Fase Quântica, em inglês {\itshape Quantum Phase Estimation} (QPE), permite estimar a fase aplicada por um operador desconhecido. Ela se baseia em incorporar portas não controladas ao algoritmo da Transformada Quântica de Fourier \cite{coppersmith1994qft}, que, por sua vez, utiliza algumas portas controladas. \cite{kitaev1995qpe}

\section{Execução dos Experimentos}\label{sec:execution}

Para evitar flutuações no tempo de execução, a ordem de cada simulação foi organizada para que os simuladores sejam testados sequencialmente. Ou seja, primeiro foram feitas todas as simulações do QMDD, depois todas do simulador denso e por fim todas do simulador esparso. Assim, questões práticas da implementação de cada simulador em relação à arquitetura do computador utilizado terão menos interferência um no outro. Por exemplo, um simulador que utiliza a memória de forma mais organizada não vai ser afetado por um simulador que deixa a memória do computador fragmentada após sua execução.

Além disso, todas as instâncias de simulação foram realizadas múltiplas vezes, utilizando rodadas de aquecimento para reduzir picos nas rodadas iniciais. O número de execuções foi variado de acordo com a quantia de qubits na simulação: quanto mais qubits, menos rodadas. Isso pois, dado o maior número de qubits, o tempo de execução se torna maior, fazendo com que a variação entre execuções seja menos relevante.

Desse modo, cada combinação de simulador e algoritmo foi executada usando diferentes números de qubits no algoritmo (no FALQON, o número de camadas foi fixado em 10). Para isso, cada combinação iniciou seus experimentos com uma quantia mínima de qubits (1 para Grover, 3 para FALQON e 2 para QPE) e, após cada execução bem sucedida, o número de qubits era incrementado em um. Uma execução falha caso o tempo de execução ultrasse um tempo limite previamente estabelecido.

Além disso, também foi feita a comparação da execução do FALQON ao variar o número de iterações utilizadas no processo de otimização. Nesse caso, o aumento da demora é linear em relação ao número de camadas ---com o dobro de camadas, espera-se que a simulação demore o dobro do tempo--- enquanto mantém-se um número constante de qubits. Para esse teste, o número de qubits foi fixado em 18 por gerar um tempo base grande o suficiente para ser estável mas ainda pequeno o suficiente para manter conveniente as múltiplas execuções com diferentes números de camadas.

\section{Resultados}\label{sec:results}

Nesta seção, os tempos apresentados correspondem à média das execuções registradas para cada combinação de algoritmo, simulador e tamanho do circuito. Quando uma combinação foi executada mais de uma vez, as médias das execuções foram agrupadas novamente pelo número de qubits ou pelo número de camadas. Experimentos que ultrapassaram o limite estabelecido não foram incluídos em intervalos maiores.

\subsection{Variação do número de qubits}

\subsubsection{Grover}

Na Figura~\ref{fig:grover-by-qubits}, observa-se que o simulador QMDD apresentou resultado competitivo com os demais simuladores em circuitos de até 14 qubits. Nesse ponto, o tempo médio foi de aproximadamente $0{,}49$~s no QMDD, contra $0{,}31$~s no simulador denso e $3{,}21$~s no esparso. Até 13 qubits, o simulador QMDD permaneceu próximo ao denso e foi mais rápido que o esparso.

A partir de 15 qubits, o crescimento da quantidade de qubits de controle produziu um aumento expressivo no tempo do QMDD. Com 18 qubits, o tempo médio foi de $42{,}85$~s no QMDD, enquanto os simuladores denso e esparso levaram, respectivamente, $2{,}72$~s e $349{,}87$~s. O QMDD continuou sendo executado com 19 qubits, em $92{,}09$~s, mas não concluiu a combinação com 20 qubits dentro do limite definido.

Nesse cenário, o simulador denso alcançou o maior intervalo, tendo sido executado até 23 qubits, com tempo médio de $890{,}13$~s nessa dimensão. O simulador esparso foi executado até 19 qubits, atingindo $1174{,}49$~s nesse último tamanho concluído. Assim, o crescimento dos controles do algoritmo de Grover foi particularmente desfavorável à implementação atual do QMDD.

\begin{figure}[ht]
  \centering
  \includegraphics[width=0.88\linewidth]{images/grover_by_qubits.png}
  \caption{Tempo médio do algoritmo de Grover em função do número de qubits.
  O eixo vertical utiliza escala logarítmica.}
  \label{fig:grover-by-qubits}
\end{figure}

\subsubsection{FALQON}

O comportamento obtido com o FALQON foi mais favorável ao QMDD do que o observado no algoritmo de Grover, como ilustrado pela Figura~\ref{fig:falqon-by-qubits}. Na faixa inicial, os três simuladores apresentaram tempos próximos. Com 14 qubits, os tempos médios foram aproximadamente $0{,}32$~s para o QMDD, $0{,}32$~s para o denso e $0{,}82$~s para o esparso.

Com 17 qubits, o QMDD levou $1{,}81$~s, contra $2{,}12$~s do denso e $6{,}40$~s do esparso. Com 18 qubits, o tempo médio foi de $2{,}98$~s no QMDD, $3{,}08$~s no denso e $11{,}59$~s no esparso. A partir de 19 qubits, o denso novamente apresentou vantagem numérica sobre o QMDD. Mesmo assim, o método QMDD continuou substancialmente mais rápido que o esparso até as maiores dimensões em que ambos concluíram o experimento.

O simulador QMDD foi concluído até 25 qubits, com tempo médio de $309{,}56$~s. Nessa dimensão, o denso levou $230{,}26$~s e o esparso não concluiu a execução dentro do limite. O simulador denso chegou a 27 qubits, atingindo $1362{,}15$~s, enquanto o esparso foi executado até 24 qubits, com $1475{,}26$~s. Esses resultados indicam que o FALQON permite ao QMDD compartilhar partes repetidas da estrutura com maior eficácia do que o Grover.

\begin{figure}[ht]
  \centering
  \includegraphics[width=0.88\linewidth]{images/falqon_by_qubits.png}
  \caption{Tempo médio do algoritmo FALQON em função do número de qubits.
  O eixo vertical utiliza escala logarítmica.}
  \label{fig:falqon-by-qubits}
\end{figure}

\subsubsection{Estimação de fase quântica}

Na Estimação de Fase Quântica, os três simuladores apresentaram desempenhos próximos até aproximadamente 16 qubits, como mostrado na Figura~\ref{fig:qpe-by-qubits}. Com 18 qubits, o simulador esparso levou $0{,}49$~s, o denso $0{,}15$~s e o QMDD $7{,}90$~s. Nessa faixa intermediária, o denso apresentou o melhor crescimento. Com 20 qubits, o denso levou $0{,}17$~s, o esparso $1{,}96$~s e o QMDD $115{,}61$~s.

A partir desse ponto, o tempo do QMDD passou a crescer rapidamente e não concluiu a combinação com 21 qubits. Em contraste, o simulador esparso chegou a 27 qubits, com tempo médio de $506{,}62$~s, e o denso alcançou 30 qubits, com $135{,}80$~s. Assim, a combinação de portas controladas com operações da Transformada Quântica de Fourier foi desfavorável ao QMDD no intervalo testado, embora o método permanecesse próximo aos demais em circuitos menores.

\begin{figure}[ht]
  \centering
  \includegraphics[width=0.88\linewidth]{images/qpe_by_qubits.png}
  \caption{Tempo médio da Estimação de Fase Quântica em função do número de qubits.
  O eixo vertical utiliza escala logarítmica.}
  \label{fig:qpe-by-qubits}
\end{figure}

\subsection{Variação do número de camadas no FALQON}

No experimento de variação das camadas, o número de qubits foi mantido em 18 e o número de iterações do FALQON foi aumentado de 100 a 5400. Cada aumento de 100 camadas repetiu uma estrutura já construída nas iterações anteriores, de modo que o crescimento observado deveu-se principalmente ao número de aplicações do circuito e não ao aumento do vetor de estado.

O tempo do QMDD apresentou crescimento aproximadamente linear, como pode ser visto na Figura~\ref{fig:falqon-by-layers}. Com 100 camadas, foram necessários $29{,}67$~s; com 1000 camadas, $297{,}89$~s; com 3000 camadas, $1107{,}21$~s; e com 5000 camadas, $2181{,}50$~s. Isso corresponde a uma razão próxima de 0{,}3~s por camada no ponto intermediário do experimento.

Na comparação com o simulador denso, o QMDD levou menos tempo a partir de aproximadamente 1500 camadas. Com 1000 camadas, o denso levou $325{,}79$~s contra $297{,}89$~s do QMDD; com 1500 camadas, os tempos foram $521{,}17$~s e $483{,}68$~s, respectivamente. Com 3500 camadas, o denso alcançou $1791{,}44$~s, enquanto o QMDD levou $1313{,}37$~s.

Como observado nos resultados do experimento de variação do número de qubits para o FALQON, o simulador esparso apresentou desempenho significativamente pior que o simulador denso e o QMDD. Por isso, os resultados do simulador esparso para execuções maiores do FALQON não foram incluídos.

Em ambos os pontos, o QMDD apresentou menor tempo. Os resultados mostram que, embora o custo por camada seja significativo, a reutilização da estrutura no QMDD torna o aumento das iterações mais eficiente que a execução repetida pelos vetores denso e esparso no intervalo observado.

\begin{figure}[ht]
  \centering
  \includegraphics[width=0.88\linewidth]{images/falqon_by_layers.png}
  \caption{Tempo médio do FALQON com 18 qubits em função do número de camadas.
  Os pontos esparsos disponíveis no intervalo de 1000 a 1500 camadas são
  comparados separadamente no texto.}
  \label{fig:falqon-by-layers}
\end{figure}

\subsection{Síntese dos resultados}

Os resultados mostram que o desempenho do QMDD está diretamente relacionado à estrutura do circuito. No Grover e na Estimação de Fase Quântica, o crescimento dos qubits de controle fez o número de vértices crescer rapidamente e reduziu a competitividade do método. No FALQON, em que o diagrama construído nas primeiras iterações é reutilizado nas iterações seguintes, o simulador manteve melhor desempenho e apresentou vantagem sobre os demais em partes relevantes do intervalo.

O simulador denso alcançou o maior número de qubits no Grover e na Estimação de Fase Quântica. Isso é coerente com sua representação direta do vetor de estado e com a possibilidade de distribuir independentemente as operações sobre suas posições em arquiteturas com vários núcleos. O simulador esparso apresentou vantagem em algumas dimensões da Estimação de Fase Quântica, mas seu custo cresceu de forma mais intensa no Grover.

Portanto, o novo simulador não supera os simuladores existentes em todos os circuitos. Sua eficiência foi maior em problemas em que padrões estruturais podem ser reutilizados e em que o número de controles não domina o crescimento do diagrama. O FALQON e o experimento com 18 qubits variando o número de camadas representaram os cenários mais favoráveis à implementação atual.

